% =============================================================================
\section{Detalles de Implementación}
% =============================================================================

La simulación fue implementada íntegramente en Python~3.13, sin el uso de
librerías de simulación externas (p.ej.\ SimPy), con el objetivo de
exhibir explícitamente cada mecanismo del paradigma de eventos discretos.
El código fuente se organiza en seis paquetes independientes:
\texttt{rng} (generación aleatoria), \texttt{engine} (motor DES),
\texttt{model} (dominio del restaurante), \texttt{metrics} (colección
de estadísticos), \texttt{experiments} (ejecución de réplicas) y
\texttt{analysis} (estadística y visualización).

% ---------------------------------------------------------------------------
\subsection{Generador de números aleatorios (LCG)}
% ---------------------------------------------------------------------------

La base del sistema estocástico es un \textbf{Generador Congruencial Lineal}
(LCG), definido por la recurrencia~\citep{knuth1997art}:
\begin{equation}
  X_{n+1} = (A \cdot X_n + C) \bmod M, \qquad U_n = \frac{X_n}{M},
  \label{eq:lcg}
\end{equation}
con los parámetros de Knuth usados también por la biblioteca estándar de C
en Linux (glibc):
\[
  A = 1\,664\,525, \quad C = 1\,013\,904\,223, \quad M = 2^{32}.
\]
Estos parámetros satisfacen los criterios de Hull--Dobell para período
máximo igual a $M = 4\,294\,967\,296$, garantizando que la secuencia no se
repite dentro de ninguna simulación práctica.

% ---------------------------------------------------------------------------
\subsection{Método de la transformada inversa}
% ---------------------------------------------------------------------------

Todas las variables aleatorias se generan mediante el método de la
transformada inversa, cuya validez se sustenta en la siguiente proposición
\citep{curso_conf2,ross2013simulation}.

\begin{quote}
\textbf{Proposición (Transformada Inversa).}
Sea $U \sim \mathcal{U}(0,1)$.
Para cualquier función de distribución acumulada continua $F$, la variable
aleatoria $X = F^{-1}(U)$ tiene distribución $F$.
\end{quote}

\textit{Demostración.}
Sea $F_X(x) = P[X \le x]$.
Entonces:
\[
  F_X(x) = P\!\left[F^{-1}(U) \le x\right]
          = P\!\left[U \le F(x)\right]
          = F(x), \qquad \square
\]
donde la última igualdad es consecuencia de que $U$ es uniforme en $(0,1)$.

\paragraph{Distribución exponencial.}
Para modelar los tiempos entre llegadas se emplea
$X \sim \operatorname{Exp}(\lambda)$, con función de distribución
$F(x) = 1 - e^{-\lambda x}$.
Invirtiendo:
\begin{equation}
  X = -\frac{1}{\lambda}\ln(U).
  \label{eq:exp}
\end{equation}

\paragraph{Distribución uniforme.}
Los tiempos de preparación siguen
$X \sim \mathcal{U}[a,b]$, con $F(x) = (x-a)/(b-a)$.
La inversión da:
\begin{equation}
  X = a + (b - a)\,U.
  \label{eq:unif}
\end{equation}

% ---------------------------------------------------------------------------
\subsection{Streams independientes y Common Random Numbers (CRN)}
% ---------------------------------------------------------------------------

Para garantizar comparaciones estadísticamente válidas entre escenarios se
implementa la técnica de \textbf{números aleatorios comunes}
(CRN)~\citep{law2015simulation}.
Se mantienen tres streams LCG independientes:
\begin{itemize}
  \item \texttt{arrivals}: genera los tiempos entre llegadas sucesivas,
  \item \texttt{type}: determina si cada cliente pide sándwich o sushi,
  \item \texttt{service}: genera la duración de preparación de cada pedido.
\end{itemize}

Cada stream parte de una semilla separada por $10^6$ posiciones:
\[
  \texttt{seed}_{\text{arrivals}} = s, \quad
  \texttt{seed}_{\text{type}}     = s + 10^6, \quad
  \texttt{seed}_{\text{service}}  = s + 2 \times 10^6.
\]
Con ${\sim}200$ clientes por día y 3 streams, el consumo máximo por réplica
es de ${\approx}600$ números, muy inferior al gap de $10^6$.

El invariante clave del CRN es que el tiempo de servicio de cada cliente se
sortea \emph{en el momento de su llegada}, no cuando inicia el servicio.
Así, el cliente $k$ siempre consume el $k$-ésimo número del stream
\texttt{service} independientemente del escenario.
Esto induce correlación positiva $\rho > 0$ entre las réplicas pareadas de
A y B, reduciendo la varianza del estimador de la diferencia:
\[
  \mathbb{V}(\bar{D}) = \mathbb{V}(\bar{A}) + \mathbb{V}(\bar{B}) - 2\,\operatorname{Cov}(A,B).
\]

% ---------------------------------------------------------------------------
\subsection{Tipos de eventos}
% ---------------------------------------------------------------------------

El modelo define cinco tipos de eventos, descritos en \cref{tab:eventos}.

\begin{table}[H]
  \centering
  \caption{Tipos de eventos del modelo DES de La Cocina de Kojo.}
  \label{tab:eventos}
  \begin{tabular}{lp{9cm}}
    \toprule
    \textbf{Tipo} & \textbf{Descripción y acciones principales} \\
    \midrule
    \textsc{llegada}        & Arribo de un nuevo cliente. Se determina su tipo
                              (Bernoulli) y su tiempo de servicio (Uniforme, al llegar).
                              Si hay servidor libre se asigna; si no, el cliente
                              entra a la cola FIFO.
                              Se programa la siguiente llegada. \\
    \addlinespace
    \textsc{partida}        & Fin del servicio de un cliente. El servidor queda
                              libre. Si la cola no está vacía se atiende al
                              siguiente cliente. Si el empleado está marcado
                              para retiro, se elimina. \\
    \addlinespace
    \textsc{inicio\_pico}   & Comienzo de una hora pico. En el Escenario B se
                              añade el empleado extra; si hay cola, comienza a
                              atender de inmediato. \\
    \addlinespace
    \textsc{fin\_pico}      & Fin de una hora pico. En el Escenario B el empleado
                              extra es marcado para retiro si está ocupado, o
                              retirado inmediatamente si está libre. \\
    \addlinespace
    \textsc{fin\_simulacion}& Cierre del día (t = 660 min). Los servicios en curso
                              se completan; los clientes en cola no reciben servicio. \\
    \bottomrule
  \end{tabular}
\end{table}

% ---------------------------------------------------------------------------
\subsection{Algoritmo del bucle DES}
% ---------------------------------------------------------------------------

El \cref{alg:des} muestra el bucle principal de simulación, fiel al
algoritmo de avance de tiempo por próximo evento descrito
en~\citet{banks2010discrete}.

\begin{algorithm}[H]
\caption{Bucle principal de la simulación DES}
\label{alg:des}
\begin{algorithmic}[1]
\Require configuración $\mathtt{cfg}$, streams $\mathtt{rng}$
\Ensure \texttt{ReplicationResult} con métricas del día
\State Crear empleados base; programar primera \textsc{llegada} y eventos de pico
\While{la lista de eventos futuros (FEL) no está vacía}
  \State $e \gets$ extraer evento de mínimo tiempo de la FEL
  \State Avanzar el reloj: $t \gets e.\mathtt{time}$
  \If{$e$ es \textsc{llegada}}
    \State Sortear tipo de cliente y $s \gets$ tiempo de servicio (Transformada Inversa)
    \State Crear cliente con $\mathtt{arrival\_time} = t$, $\mathtt{service\_time} = s$
    \If{existe servidor libre}
      \State Asignar cliente al servidor; programar \textsc{partida} en $t + s$
    \Else
      \State Encolar cliente; registrar cambio de cola
    \EndIf
    \State Programar siguiente \textsc{llegada} en $t + \operatorname{Exp}(\lambda(t))$
  \ElsIf{$e$ es \textsc{partida}}
    \State Liberar servidor; registrar $\mathtt{departure\_time} = t$
    \If{cola no vacía \textbf{y} $t \le T$}
      \State Desencolar siguiente cliente; asignarlo; programar \textsc{partida}
    \EndIf
  \ElsIf{$e$ es \textsc{inicio\_pico} \textbf{y} $\mathtt{extra\_staff} > 0$}
    \State Crear empleado extra; si hay cola, asignarle cliente inmediatamente
  \ElsIf{$e$ es \textsc{fin\_pico} \textbf{y} $\mathtt{extra\_staff} > 0$}
    \State Marcar empleado extra para retiro (o retirar si está libre)
  \EndIf
\EndWhile
\State \Return $\mathtt{MetricsCollector.summarize}()$
\end{algorithmic}
\end{algorithm}

% ---------------------------------------------------------------------------
\subsection{Proceso de Poisson no homogéneo}
% ---------------------------------------------------------------------------

El proceso de llegadas es un proceso de Poisson no homogéneo con función de
intensidad por tramos constantes \citep{curso_conf2}:
\[
  \lambda(t) = \begin{cases}
    0{,}30 & \text{si } t \in [90,210) \cup [420,540), \\
    0{,}17 & \text{en otro caso.}
  \end{cases}
\]
En la implementación, cada evento \textsc{llegada} programa el siguiente
usando $\lambda(t_{\text{actual}})$: en cuanto el reloj supera una frontera
de pico, la siguiente llegada ya utiliza la tasa correcta.

% ---------------------------------------------------------------------------
\subsection{Cálculo de la longitud media de cola}
% ---------------------------------------------------------------------------

La longitud media de cola se estima como el cociente entre el área bajo la
curva $Q(t)$ (número de clientes en cola en el instante $t$) y la duración
del día $T = 660$ min:
\begin{equation}
  \bar{L}_q = \frac{1}{T}\int_0^T Q(t)\,dt
  \approx \frac{1}{T}\sum_k Q(t_k)\,(t_{k+1} - t_k),
  \label{eq:area_cola}
\end{equation}
donde la suma es una integral de Riemann por rectángulos actualizada en cada
evento de cambio de cola.
Este estimador es equivalente al que se obtiene aplicando la Ley de
Little~\citep{little1961proof}: $L_q = \lambda\,W_q$.

% ---------------------------------------------------------------------------
\subsection{Protocolo de réplicas y parada}
% ---------------------------------------------------------------------------

Se ejecutan $n = 30$ réplicas independientes para cada escenario.
La semilla de la réplica $i$ es $s_i = 42 + i \times 10^5$, garantizando
una separación mayor que el consumo por réplica (${\approx}600$ números).
El criterio de parada formal propuesto en~\citet{curso_temas} (cap.~4) es:
\begin{equation}
  \frac{S}{\sqrt{n}} < d,
  \label{eq:parada_temas}
\end{equation}
donde $d$ es la desviación estándar máxima admisible para el estimador.
Este criterio se analiza en detalle en la \cref{sec:parada}.
